\section{Introduction}
\label{sec:introduction}

The \texttt{penguins} dataset ships with the \texttt{seaborn} library and contains morphological measurements of 344 adult penguins from three species (\textit{Adelie}, \textit{Chinstrap} and \textit{Gentoo}) collected on three islands of the Palmer Archipelago. For each individual the dataset records four continuous physical traits (bill length, bill depth, flipper length and body mass) together with two categorical attributes: island of origin and sex. Species are strongly separated by morphology, which makes this dataset a common benchmark for classification. Males and females also differ systematically in size, which raises the question of how much a model gains from knowing sex.

This report answers two questions: the first is how well species can be predicted from physical traits alone and how much that changes when sex is added; the second is whether the same pattern holds when body mass is the regression target instead. Four supervised models are compared: logistic regression and a decision tree for classification, and linear regression and a random forest for regression. A repeated A/B design is used to check whether the conclusions depend on which rows end up in the training set.

The rest of the report is organized as follows. Section~\ref{sec:data_loading} covers data selection, loading and preprocessing. Section~\ref{sec:classification} and Section~\ref{sec:regression} present the classification and regression studies. Section~\ref{sec:robustness} checks whether the conclusions hold when the training set is split into two disjoint halves. Hyperparameter optimization of the decision tree is covered in Section~\ref{sec:hyperparameter_tuning}. Section~\ref{sec:conclusion} closes with the main findings.